\chapter{Формулы Кардано и Феррари}\label{chap:c23-radicals}\markright{Глава 39}
\section{Общие формулы решения в радикалах}

\par\addvspace{5pt}\begin{lrbox}{\CatFlowBox}\begin{minipage}[t]{\linewidth}
\setlength{\parskip}{5pt}\setlength{\columnsep}{4.5mm}\setlength{\intextsep}{0pt}
\begin{wrapfigure}[12]{r}{54mm}
\includegraphics[width=54mm,height=58mm,keepaspectratio]{images/unique/art-79.png}
\end{wrapfigure}
Общая формула решения должна работать не только тогда, когда корень удаётся угадать или многочлен уже удобно разложен на множители. Она должна начинаться с коэффициентов, указывать допустимые вычисления и возвращать все решения. После введения комплексных чисел можно последовательно извлекать квадратные и кубические корни даже из отрицательных действительных и из недействительных чисел. Однако у такого извлечения обычно несколько значений. Поэтому существенная часть формулы --- правило согласования выбранных значений, а не только привычная запись с радикалами.

При ручном решении полезно сначала проверить простые возможности: свободный член может быть нулевым, многочлен может содержать только чётные степени или иметь удобный рациональный корень. Общие формулы нужны, когда эти сокращения не помогают, и одновременно объясняют, почему степени три и четыре разрешимы через радикалы. Нахождение одного решения ещё не завершает задачу: нужно восстановить сдвиг переменной, перечислить остальные решения и проверить, не совпали ли некоторые из них. Ниже полнота получается из явного разложения исходного многочлена на линейные или квадратные множители.
\par\WFclear\end{minipage}\end{lrbox}%
\ifdim\dimexpr\ht\CatFlowBox+\dp\CatFlowBox\relax<64mm\dp\CatFlowBox=\dimexpr64mm-\ht\CatFlowBox\relax\fi%
\noindent\usebox{\CatFlowBox}\par\addvspace{5pt}
\begin{definition}
\term{Формула в радикалах} (formula by radicals) использует коэффициенты уравнения, рациональные числа и конечную последовательность сложений, вычитаний, умножений, допустимых делений и извлечений корней натуральных степеней. Промежуточные числа могут быть комплексными.
\end{definition}
\begin{definition}
\term{Кратность корня} (multiplicity of a root) \(c\) ненулевого многочлена \(P\) --- натуральное число \(m\ge1\), для которого \(P(x)=(x-c)^mQ(x)\), где \(Q(c)\ne0\). При перечислении с кратностями такое значение записывают \(m\) раз.
\end{definition}

В формулах этого раздела запись выбора корня означает любое число с указанной степенью; она не означает автоматически главное значение. Все коэффициенты допускаются комплексными. Если нужны только действительные решения, их отбирают из полного комплексного списка.

\section{Приведение кубического уравнения}

Для \(ax^3+bx^2+cx+d=0\), \(a\ne0\), сначала разделим на \(a\) и обозначим
\[
A=\frac ba,\qquad B=\frac ca,\qquad C=\frac da.
\]
\begin{theorem}[Устранение квадратного слагаемого]\label{thm:legacy-c23-10}
Замена \(x=y-\frac A3\) превращает \(x^3+Ax^2+Bx+C=0\) в
\[
y^3+py+q=0,\qquad
p=B-\frac{A^2}{3},\qquad q=\frac{2A^3}{27}-\frac{AB}{3}+C.
\]
\end{theorem}
\begin{proof}
После подстановки коэффициент при \(y^2\) равен \(-A+A=0\). Коэффициент при \(y\) равен \(\frac{A^2}{3}-\frac{2A^2}{3}+B=p\). Свободный член равен \(-\frac{A^3}{27}+\frac{A^3}{9}-\frac{AB}{3}+C=q\). Коэффициент при кубе остаётся единицей.
\end{proof}

\begin{definition}
Кубическое уравнение без квадратного слагаемого называют \term{приведённым к неполной форме} (depressed cubic). Для него обозначим
\[
\Delta=\left(\frac q2\right)^2+\left(\frac p3\right)^3,
\qquad \omega=-\frac12+\frac{\sqrt3}{2}i.
\]
Здесь \(\omega^3=1\), \(\omega\ne1\), \(1+\omega+\omega^2=0\). Величина \(\Delta\) отличается от обычного дискриминанта кубика: последний равен \(-108\Delta\).
\end{definition}

\section{Формула Кардано и согласование кубических корней}

\begin{theorem}[Кардано; Cardano's formula]\label{thm:legacy-c23-11}
Для \(y^3+py+q=0\) полный алгоритм таков.

Если \(p=0\), решаем \(y^3=-q\). При \(q=0\) единственное значение решения --- \(0\), кратность три. При \(q\ne0\) выбираем любое \(u\) с \(u^3=-q\); все решения --- \(u,\omega u,\omega^2u\).

Если \(p\ne0\), выбираем любое \(s\) с \(s^2=\Delta\), затем любое \(u\), для которого
\[
u^3=-\frac q2+s,
\qquad v=-\frac p{3u}.
\]
Тогда все три корня с учётом кратностей имеют вид
\[
y_k=\omega^ku+\omega^{-k}v,\qquad k=0,1,2.
\]
Все корни первоначального уравнения равны \(x_k=y_k-\frac A3\).
\end{theorem}
\begin{proof}
Ветка \(p=0\) уже решена формулой всех комплексных кубических корней. Пусть \(p\ne0\). Обозначим \(U=-\frac q2+s\), \(V=-\frac q2-s\). Тогда
\[
U+V=-q,\qquad UV=\frac{q^2}{4}-\Delta=-\frac{p^3}{27}\ne0.
\]
Поэтому \(U\ne0\), любой его кубический корень \(u\) ненулевой, и деление при выборе \(v\) законно. Получаем \(uv=-\frac p3\), \(v^3=V\), \(u^3+v^3=-q\). Для любого согласованного выбора
\[
(u+v)^3+p(u+v)+q
=u^3+v^3+q+(3uv+p)(u+v)=0.
\]
Умножение \(u\) на \(\omega^k\), а \(v\) на \(\omega^{-k}\) сохраняет эти условия. Для полноты используется тождество
\[
\prod_{k=0}^{2}\bigl(t-\omega^ku-\omega^{-k}v\bigr)
=t^3-3uv\,t-(u^3+v^3)=t^3+pt+q.
\]
Оно получается раскрытием трёх множителей с учётом \(1+\omega+\omega^2=0\). Произведение равно нулю только при нулевом множителе, так что других решений нет. Само тождество также сохраняет кратности.
\end{proof}

\par\addvspace{5pt}\begin{lrbox}{\CatFlowBox}\begin{minipage}[t]{\linewidth}
\setlength{\parskip}{5pt}\setlength{\columnsep}{4.5mm}\setlength{\intextsep}{0pt}
\begin{wrapfigure}[11]{r}{54mm}
\includegraphics[width=54mm,height=58mm,keepaspectratio]{images/unique/art-100.png}
\end{wrapfigure}
Нельзя независимо взять любые кубические корни из \(-\frac q2+s\) и \(-\frac q2-s\): их произведение обязано быть \(-\frac p3\). Выбор \(v\) через \(u\) обеспечивает это автоматически. Замена \(u\) на другой кубический корень лишь переставляет три полученных значения; смена знака \(s\) также не меняет полный набор, поскольку доказанное произведение равно тому же многочлену.

При \(\Delta=0\), \(p\ne0\), можно сразу показать кратности. Положим \(c=-\frac{3q}{2p}\). Из \(q^2=-\frac{4p^3}{27}\) следуют \(p=-3c^2\), \(q=2c^3\), поэтому
\[
y^3+py+q=(y-c)^2(y+2c).
\]
Корень \(c\) двойной, \(-2c\) простой. Здесь \(c\ne0\); случай \(p=q=0\) разобран отдельно. Для действительных коэффициентов даже при действительных ответах промежуточные радикалы могут оказаться комплексными: это не ошибка и не основание отбрасывать решения.
\par\WFclear\end{minipage}\end{lrbox}%
\ifdim\dimexpr\ht\CatFlowBox+\dp\CatFlowBox\relax<64mm\dp\CatFlowBox=\dimexpr64mm-\ht\CatFlowBox\relax\fi%
\noindent\usebox{\CatFlowBox}\par\addvspace{5pt}
\section{Приведение уравнения четвёртой степени}

Для \(ax^4+bx^3+cx^2+dx+e=0\), \(a\ne0\), обозначим
\[
A=\frac ba,\quad B=\frac ca,\quad C=\frac da,\quad D=\frac ea.
\]
\begin{theorem}[Устранение кубического слагаемого]\label{thm:legacy-c23-12}
Замена \(x=y-\frac A4\) превращает нормированное уравнение в
\[
y^4+py^2+qy+r=0,
\]
где
\[
p=B-\frac{3A^2}{8},\qquad q=C-\frac{AB}{2}+\frac{A^3}{8},
\]
\[
r=D-\frac{AC}{4}+\frac{A^2B}{16}-\frac{3A^4}{256}.
\]
\end{theorem}
\begin{proof}
После раскрытия степеней коэффициент при \(y^3\) равен \(-A+A=0\). При \(y^2\) получаем \(\frac{6A^2}{16}-\frac{3A^2}{4}+B=p\); при \(y\) --- \(-\frac{A^3}{16}+\frac{3A^3}{16}-\frac{AB}{2}+C=q\). Свободный член равен \(\frac{A^4}{256}-\frac{A^4}{64}+\frac{A^2B}{16}-\frac{AC}{4}+D=r\).
\end{proof}

Сначала проверим \(q=0\). Тогда уравнение биквадратное: \(T=y^2\) удовлетворяет \(T^2+pT+r=0\). Выберем любой \(d_0\) с \(d_0^2=p^2-4r\) и положим
\[
T_1=\frac{-p+d_0}{2},\qquad T_2=\frac{-p-d_0}{2}.
\]
Все решения --- два квадратных корня каждого \(T_j\), то есть \(\pm\sqrt{T_1}\), \(\pm\sqrt{T_2}\) с пониманием выбора комплексных корней. Полнота следует из \((y^2-T_1)(y^2-T_2)=0\). Если \(T_j=0\), оба знака дают одно значение \(0\), которое учитывается дважды; если \(T_1=T_2\), кратности также сохраняются.

\section{Метод Феррари: от четвёртой степени к кубической}

\par\addvspace{5pt}\begin{lrbox}{\CatFlowBox}\begin{minipage}[t]{\linewidth}
\setlength{\parskip}{5pt}\setlength{\columnsep}{4.5mm}\setlength{\intextsep}{0pt}
\begin{wrapfigure}[11]{l}{54mm}
\includegraphics[width=54mm,height=58mm,keepaspectratio]{images/unique/art-80.png}
\end{wrapfigure}
При решении четвёртой степени удобно искать не сразу четыре линейных множителя, а два квадратных. Исчезнувшее кубическое слагаемое подсказывает, что линейные коэффициенты этих множителей должны быть противоположными. Остаётся подобрать их величину и свободные члены так, чтобы совпали ещё три коэффициента. Это приводит к вспомогательному кубическому уравнению. Его решения уже можно получить формулой Кардано, после чего остаются только квадратные уравнения.

Выбор одного вспомогательного корня не означает выбора лишь части ответов исходной задачи. Он задаёт один способ разложить тот же многочлен на два квадратных множителя, а каждый из них решается полностью. Разные допустимые выборы могут менять промежуточные числа и порядок ответов, но проверяемое произведение остаётся прежним. Поэтому важно выписать оба квадратных множителя, взять по два корня каждого и лишь затем объединить совпадающие значения с учётом их кратностей. Отдельная биквадратная ветка выше нужна ещё и потому, что в общем преобразовании появится деление на выбранный параметр; условие его ненулевости должно быть доказано до деления.
\par\WFclear\end{minipage}\end{lrbox}%
\ifdim\dimexpr\ht\CatFlowBox+\dp\CatFlowBox\relax<64mm\dp\CatFlowBox=\dimexpr64mm-\ht\CatFlowBox\relax\fi%
\noindent\usebox{\CatFlowBox}\par\addvspace{5pt}
\begin{definition}
\term{Кубическая резольвента} (resolvent cubic) для \(y^4+py^2+qy+r\) в используемой параметризации --- уравнение
\[
z^3+2pz^2+(p^2-4r)z-q^2=0.
\]
Его корень будет равен квадрату линейного коэффициента в квадратных множителях.
\end{definition}

\begin{theorem}[Феррари; Ferrari's formula]\label{thm:legacy-c23-13}
Пусть \(q\ne0\). Выберем любой корень \(z\) указанной резольвенты и любой \(s\) с \(s^2=z\). Тогда \(z\ne0\), \(s\ne0\). Положим
\[
t=\frac{p+z-\frac qs}{2},\qquad
u=\frac{p+z+\frac qs}{2}.
\]
Выполняется разложение
\[
y^4+py^2+qy+r=(y^2+sy+t)(y^2-sy+u).
\]
Выберем любые \(d_1,d_2\) с
\[
d_1^2=-2p-z+\frac{2q}{s},\qquad
d_2^2=-2p-z-\frac{2q}{s}.
\]
Все четыре корня с учётом кратностей:
\[
y_1=\frac{-s+d_1}{2},\quad y_2=\frac{-s-d_1}{2},
\]
\[
y_3=\frac{s+d_2}{2},\quad y_4=\frac{s-d_2}{2}.
\]
Для исходного уравнения к каждому значению применяют \(x=y-\frac A4\).
\end{theorem}
\begin{proof}
Значение резольвенты в нуле равно \(-q^2\ne0\), поэтому ни один её корень не нулевой. По определению \(s\ne0\), и все деления допустимы. Разложение проверяется четырьмя коэффициентами:
\[
t+u-s^2=p,\qquad s(u-t)=q,
\]
\[
tu=\frac{(p+z)^2-\frac{q^2}{z}}4=r.
\]
Последнее равенство есть резольвента, разделённая на \(4z\). Коэффициент при \(y^3\) равен нулю, при \(y^4\) --- единице. Значит, произведение действительно совпадает с исходным многочленом. Квадратное уравнение \(y^2+sy+t=0\) имеет дискриминант \(s^2-4t=d_1^2\) и даёт первые два корня. Для \(y^2-sy+u=0\) дискриминант равен \(s^2-4u=d_2^2\), получаются последние два. Свойство нулевого произведения доказывает полноту списка, в том числе при нулевом дискриминанте или общем корне двух множителей.
\end{proof}

\par\addvspace{5pt}\begin{lrbox}{\CatFlowBox}\begin{minipage}[t]{\linewidth}
\setlength{\parskip}{5pt}\setlength{\columnsep}{4.5mm}\setlength{\intextsep}{0pt}
\begin{wrapfigure}[13]{l}{54mm}
\includegraphics[width=54mm,height=58mm,keepaspectratio]{images/unique/art-101.png}
\end{wrapfigure}
Для вычисления корня резольвенты можно непосредственно применить уже доказанный кубический алгоритм. Если желательно сначала убрать её квадратное слагаемое, положим \(z=h-\frac{2p}{3}\). Получится
\[
h^3+Ph+Q=0,\qquad P=-\frac{p^2}{3}-4r,
\]
\[
Q=-\frac{2p^3}{27}+\frac{8pr}{3}-q^2.
\]
В ветке \(P=0\) используют соответствующий простой случай Кардано. Существование хотя бы одного корня резольвенты следует из полного кубического алгоритма, а не предполагается дополнительно. Знак \(s\) можно выбирать любой: его смена переставляет квадратные множители. Знак каждого \(d_j\) лишь переставляет два корня своего множителя. В отличие от двух кубических радикалов Кардано, квадратные радикалы \(d_1,d_2\) здесь не требуют взаимного согласования: из каждого в любом случае берут оба противоположных значения.
\par\WFclear\end{minipage}\end{lrbox}%
\ifdim\dimexpr\ht\CatFlowBox+\dp\CatFlowBox\relax<64mm\dp\CatFlowBox=\dimexpr64mm-\ht\CatFlowBox\relax\fi%
\noindent\usebox{\CatFlowBox}\par\addvspace{5pt}
\section{Граница радикальных формул}

\begin{theorem}[Абель--Руффини; Abel--Ruffini theorem; без доказательства]\label{thm:legacy-c23-14}
Для каждого натурального \(n\ge5\) не существует общей формулы, которая выражала бы все корни произвольного уравнения степени \(n\) через его коэффициенты и рациональные числа конечным числом арифметических операций и извлечений корней натуральных степеней. Уже среди уравнений с рациональными коэффициентами существуют уравнения, не все корни которых выражаются такими операциями из рациональных чисел.
\end{theorem}

\textit{Без доказательства.} Полное обоснование этого результата требует более глубокой алгебры, прежде всего теории расширений полей и теории Галуа; здесь оно не приводится. Приведённое объяснение границы метода не является доказательством теоремы.

Утверждение относится к общей радикальной формуле, а не к каждому отдельному уравнению. Например, \(x^5=2\) решается всеми пятью комплексными корнями из \(2\), а разложенные на множители уравнения высокой степени могут сводиться к простым. Теорема не утверждает отсутствия корней и не запрещает численные алгоритмы приближённого нахождения решений. Формулы Кардано и Феррари дают точные выражения; для практических десятичных вычислений также применяют численные методы с контролем ошибки.

\Needspace{7\baselineskip}\section*{Задачи}
\begingroup\setlength{\parskip}{2pt}
\noindent\begin{minipage}[t]{\linewidth}\textbf{39.1.} Реши в \(\mathbb C\) уравнение \(x^3-6x^2+15x-14=0\), сначала устранив квадратное слагаемое.\par\end{minipage}\par
\noindent\begin{minipage}[t]{\linewidth}\textbf{39.2.} Для \(y^3+3y-4=0\) вычисли радикалы Кардано с согласованным произведением, найди один действительный корень и остальные два комплексных.\par\end{minipage}\par
\noindent\begin{minipage}[t]{\linewidth}\textbf{39.3.} Реши \(x^3+8=0\) в комплексных числах. Объясни, почему ветка \(p=0\) не требует деления на кубический радикал.\par\end{minipage}\par
\noindent\begin{minipage}[t]{\linewidth}\textbf{39.4.} Реши \(y^3-12y+16=0\), вычисли \(\Delta\) и укажи кратности.\par\end{minipage}\par
\noindent\begin{minipage}[t]{\linewidth}\textbf{39.5.} Докажи, что множество корней \(y^3+py=0\) симметрично относительно нуля. Выпиши все корни и различи случаи \(p=0\) и \(p\ne0\).\par\end{minipage}\par
\noindent\begin{minipage}[t]{\linewidth}\textbf{39.6.}\starmark{} Реши в \(\mathbb C\) биквадратное уравнение \(y^4+2y^2+5=0\). Запиши корни через действительные неотрицательные квадратные корни.\par\end{minipage}\par
\noindent\begin{minipage}[t]{\linewidth}\textbf{39.7.} Примени метод Феррари к \(y^4+2y^2+y+2=0\), проверив, что \(z=1\) подходит для резольвенты. Выпиши все корни.\par\end{minipage}\par
\noindent\begin{minipage}[t]{\linewidth}\textbf{39.8.}\starmark{} Для \(y^4-3y^2+2y=0\) примени Феррари сначала с \(z=1\), затем с \(z=4\). Сравни наборы корней и укажи кратности.\par\end{minipage}\par
\noindent\begin{minipage}[t]{\linewidth}\textbf{39.9.} Докажи: если у \(P(y)=y^4+py^2+qy+r\) есть корень \(a\) кратности не менее трёх, то \(P(y)=(y-a)^3(y+3a)\). Найди \(p,q,r\) через \(a\) и определи, когда кратность становится четырьмя.\par\end{minipage}\par
\noindent\begin{minipage}[t]{\linewidth}\textbf{39.10.}\starmark{} Реши \(y^3-3y+1=0\), используя согласованные радикалы Кардано и тригонометрическую форму. Почему комплексные промежуточные значения не мешают получить три действительных корня?\par\end{minipage}\par
\endgroup
